%\todo{In charge: Eric, Charles, Hao. Final pass due May 18 5PM.}
The idea of hierarchical feature learning has been very successful. Among all the learning models, convolutional neural network \cite{krizhevsky2012imagenet,simonyan2014very,he2016deep} is one of the most prominent ones.
%, which learns multiple levels of abstraction of images using weight shared convolution kernels.
However, convolution does not apply to unordered point sets with distance metrics, which is the focus of our work.

A few very recent works \cite{qi2016pointnet,vinyals2015order} have studied how to apply deep learning to unordered sets. They ignore the underlying distance metric even if the point set does possess one. As a result, they are unable to capture local context of points and are sensitive to global set translation and normalization. In this work, we target at points sampled from a metric space and tackle these issues by explicitly considering the underlying distance metric in our design.

Point sampled from a metric space are usually noisy and with non-uniform sampling density. This affects effective point feature extraction and causes difficulty for learning. One of the key issue is to select proper scale for point feature design. Previously several approaches have been developed regarding this \cite{pauly2006point, mitra2004estimating, belton2006classification,demantke2011dimensionality, gressin2013towards,weinmann2015semantic} either in geometry processing community or photogrammetry and remote sensing community.
%Previously several approaches have been developed regarding this. \cite{pauly2006point, mitra2004estimating, belton2006classification} tries to fit curve or surfaces to the points and find optimal group around each point of interest according to hand designed rules. \cite{demantke2011dimensionality, gressin2013towards,weinmann2015semantic} focus on processing non-uniformly sampled LiDAR scans. Multiple feature scales are considered and the optimal one is selected using various heuristics.
In contrast to all these works, our approach learns to extract point features and balance multiple feature scales in an end-to-end fashion.

In 3D metric space, other than point set, there are several popular representations for deep learning, including volumetric grids \cite{qi2016volumetric,riegler2016octnet,wang2017cnn}, and geometric graphs \cite{bruna2013spectral,masci2015geodesic,yi2016syncspeccnn}. %Volumetric fields usually introduce quantization loss and artifacts on point cloud with irregular sparsity.
However, in none of these works, the problem of non-uniform sampling density has been explicitly considered.%, which is one of our focus in this work.

%In our work, we design novel neural network architectures that directly consume point sets sampled from a metric space and emphasize on handling non-uniform sampling density issue.

%In [cite graph CNN], a point cloud is firstly reconstructed into a mesh or built as an NN graph and a graph-CNN is then applied. After the connection of the graph is constructed, it is fixed by the first step pre-processing. In [cite OctNet and O-CNN], the point cloud is organized into an Octree and 3D convolutional neural networks are used for feature learning. However the octree construction introduces quantization loss and artifacts on point cloud with irregular sparsity. \emph{TODO: Octree based method on irregular density data?} In our work, we design novel neural network architectures that directly consume point sets sampled from a metric space.

%\paragraph{Hierarchical feature learning.}
%Recently with the combination of big data and deep learning we have seen success of hierarchical feature learning in many domains such as image [cite], text [cite] and audio understanding [cite]. Among all the deep models, convolutional neural network [cite] is the most prominent one. By sharing convolution kernels in each layer the network is able to learn multiple levels of abstraction of images. However, convolution does not apply to unordered point sets with distance metrics. In this work, we develop novel non-convolution based architectures with similar good properties as CNNs such as weight sharing, multi-level feature learning and translation invariance.

%\paragraph{Deep learning on point sets.}
%A few very recent works [cite PointNet, order matters, deep sets etc.] have studied how to apply deep learning on an unordered set. Usually there is a point independent embedding on each point to project it into a higher dimensional space followed by an aggregation step to combine information from all points. In that way the set is considered as a generic point set without distance metrics. There is no concept of local neighborhood of points, therefore the features learned are either single point embedding or global point set feature. Since each point is embedded independently the feature is sensitive to global set normalization and very sensitive to set translations. In this work, we fully respect the relations between points and design specific layers to capture local context of points. Our network is thus able to learn local point features of increasing receptive field and is invariant to set translations. \hao{also the variable sampling density problem} \hao{emphasize that PointNet is very robust, though.}

%\paragraph{Deep learning in 3D metric space.}
%The representation of data in 3D metric space can be sets of coordinates (point cloud), volumetric fields (occupancy or distance field), and geometric graphs. In [cite PointNet], the point cloud is considered as a generic unordered set without distance metrics or local neighborhood. In [cite graph CNN], a point cloud is firstly reconstructed into a mesh or built as an NN graph and a graph-CNN is then applied. After the connection of the graph is constructed, it is fixed by the first step pre-processing. In [cite OctNet and O-CNN], the point cloud is organized into an Octree and 3D convolutional neural networks are used for feature learning. However the octree construction introduces quantization loss and artifacts on point cloud with irregular sparsity. \emph{TODO: Octree based method on irregular density data?} In our work, we design novel neural network architectures that directly consume point sets sampled from a metric space. \hao{previous work do not consider variable density problem}.

%\paragraph{Feature design on non-uniform and noisy point cloud.}
%\hao{first need to say that scale and non-uniformity is a crucial problem for point cloud learning.}
%Previously several approaches have been developed for processing noisy 3D point cloud with irregular sampling issues. One of the key issue is to select proper scale for point feature design. Surface based approaches usually try to fit a curve or a surface to the point cloud \cite{pauly2006point}, finding the optimal group around each point of interest according to hand designed rules \cite{mitra2004estimating, belton2006classification}. These method are designed for smoothly varying surfaces and doesn't not generalize well for anistropic surfaces in scene scanning. In \cite{demantke2011dimensionality, gressin2013towards}, approaches are introduced to directly compute point features in LiDAR scans. To cope with the irregular sampling issue, a multi-scale strategy is proposed and optimal scale for each point is selected according to various criteria. Hand designed geometric features are extracted for each point within the optimal scale. In \cite{weinmann2015semantic}, the proposed method extracts hand designed geometric features for each point with respect to neighborhoods with different sizes and scales for scene segmentation. A multivariate filter-based feature selection scheme is adopted to select the feature from the optimal neighborhood. In contrast to all these works, our approach learns to extract point features and select neighborhood scales in an end-to-end fashion, making it more effective in real applications.

%\hao{scale problem and density problem are two problems.}